\chapter[\emj{⛰️}~Heap Sort (Ordenamiento por Montículo)]{\emj{⛰️}~Heap Sort (Ordenamiento por Montículo)}

\begin{dsaquees}

Ordenar sacando SIEMPRE el más grande primero 🏔️. Metes todos los números
en un montículo (heap), una estructura donde el mayor siempre está en la
cima. Sacas la cima (el máximo), la pones al final, reacomodas el heap, y
repites. Poco a poco el arreglo queda ordenado de menor a mayor.

Su gracia: garantiza $O(n \log n)$ SIEMPRE (nunca se degrada como Quick
Sort) y ordena in-place, sin memoria extra.

\end{dsaquees}

\begin{dsanino}

Imagina una torre de cajas 📦 donde la regla es que cada caja de arriba
siempre pesa más que las de abajo que cuelgan de ella. La caja más pesada
siempre queda hasta arriba.

Para ordenar: agarras la caja de arriba (la más pesada) y la pones al
final de una fila. Como quitaste la cima, subes otra caja a su lugar y
reacomodas la torre para que la regla se cumpla otra vez. Vuelves a
agarrar la nueva cima (la segunda más pesada) y la pones antes de la
anterior. Repites hasta vaciar la torre. Al final, tu fila quedó ordenada
de la más ligera a la más pesada, ¡sin necesitar otra mesa aparte!

\end{dsanino}

\begin{dsaotraforma}

Heap Sort ordena usando un \textbf{binary heap} (montículo binario)
representado en el propio arreglo. Dos fases:

\begin{enumerate}
\item \textbf{Build Max-Heap}: reorganiza el arreglo para que cumpla la
      propiedad de max-heap (cada padre $\ge$ sus hijos). Se hace en
      $O(n)$ aplicando \verb|heapify| de abajo hacia arriba.
\item \textbf{Sort}: intercambia la raíz (máximo) con el último elemento,
      reduce el tamaño del heap en 1 y aplica \verb|heapify| a la raíz.
      Repite hasta vaciar.
\end{enumerate}

En un arreglo, el nodo $i$ tiene hijos en $2i+1$ y $2i+2$, y padre en
$(i-1)/2$. No necesita punteros ni memoria extra: por eso es in-place.

\end{dsaotraforma}

\begin{dsadiagrama}
\begin{verbatim}
Array: [4, 10, 3, 5, 1]

1) Build Max-Heap:
            10                 array: [10, 5, 3, 4, 1]
          /    \
         5      3
        / \
       4   1

2) Extraer max (10) -> al final; heapify el resto:
   swap(10,1): [1, 5, 3, 4 | 10]  -> heapify -> [5, 4, 3, 1 | 10]
   swap(5,1):  [1, 4, 3 | 5, 10]  -> heapify -> [4, 1, 3 | 5, 10]
   swap(4,3):  [3, 1 | 4, 5, 10]  -> heapify -> [3, 1 | ...]
   swap(3,1):  [1 | 3, 4, 5, 10]
   Ordenado: [1, 3, 4, 5, 10]
\end{verbatim}
\end{dsadiagrama}

\begin{lstlisting}
CÓMO FUNCIONA
- heapify(i): hunde el elemento i comparándolo con sus hijos (2i+1,
  2i+2); si un hijo es mayor, intercambia y sigue bajando.
- Build: aplica heapify desde el último padre (n/2 - 1) hacia el 0.
- Sort: swap(0, fin); reduce tamaño; heapify(0). Repite.

📝 PSEUDOCÓDIGO
──────────────────────────────────────────────────────────────

función heapify(a, n, i):
    mayor = i; izq = 2*i+1; der = 2*i+2
    SI izq < n Y a[izq] > a[mayor]: mayor = izq
    SI der < n Y a[der] > a[mayor]: mayor = der
    SI mayor != i:
        intercambiar(a[i], a[mayor])
        heapify(a, n, mayor)

función heapSort(a, n):
    PARA i desde n/2-1 hasta 0: heapify(a, n, i)     // build
    PARA i desde n-1 hasta 1:
        intercambiar(a[0], a[i])                     // saca el max
        heapify(a, i, 0)                             // reacomoda

💻 CÓDIGO C++
──────────────────────────────────────────────────────────────

#include <vector>
using namespace std;

void heapify(vector<int>& a, int n, int i) {
    int mayor = i, izq = 2*i + 1, der = 2*i + 2;
    if (izq < n && a[izq] > a[mayor]) mayor = izq;
    if (der < n && a[der] > a[mayor]) mayor = der;
    if (mayor != i) { swap(a[i], a[mayor]); heapify(a, n, mayor); }
}

void heapSort(vector<int>& a) {
    int n = a.size();
    for (int i = n/2 - 1; i >= 0; i--) heapify(a, n, i);  // build max-heap
    for (int i = n - 1; i >= 1; i--) {
        swap(a[0], a[i]);         // el máximo va al final
        heapify(a, i, 0);         // reacomoda el heap reducido
    }
}

⏱️ COMPLEJIDAD (Big-O)
──────────────────────────────────────────────────────────────

  Caso        │ Tiempo        │ Nota
  ────────────┼───────────────┼──────────────────────────
  Best        │ O(n log n)    │ siempre igual
  Average     │ O(n log n)    │ sin sorpresas
  Worst       │ O(n log n)    │ NUNCA se degrada
  Espacio     │ O(1)          │ in-place
  Build heap  │ O(n)          │ (no O(n log n))

* No es estable (no conserva el orden de elementos iguales).
\end{lstlisting}

\begin{dsaentrevistas}

✅ Heap Sort vs Quick Sort: Heap Sort garantiza O(n log n) SIEMPRE, sin
   el peor caso O(n²) de Quick Sort. Pero en la práctica Quick Sort suele
   ser más rápido por mejor cache locality. Introsort usa ambos.

💡 Build-heap es O(n), no O(n log n): construir el heap desde abajo
   cuesta lineal. Es un resultado que sorprende y que preguntan. La suma
   $\sum n/2^h \cdot h$ converge a O(n).

⚠️ No es estable: si necesitas conservar el orden relativo de elementos
   iguales, Heap Sort no sirve; usa Merge Sort.

🔑 Top-K con heap: rara vez ordenas todo con Heap Sort en una entrevista;
   más común es usar un heap (priority\_queue) para los K mayores/menores
   en O(n log k). Ese es el uso real del heap.

\end{dsaentrevistas}

\begin{dsaresumen}

• Heap Sort: build max-heap, luego extrae el máximo repetidamente.
• En arreglo: hijos de i en 2i+1 y 2i+2; padre en (i-1)/2.
• O(n log n) garantizado en TODOS los casos; in-place O(1).
• Build-heap cuesta O(n) (no O(n log n)).
• No es estable.
• En entrevistas, el heap brilla más para Top-K que para ordenar todo.
════════════════════════════════════════════════════════════════

\end{dsaresumen}
